% 
% 进程包含一个线程和两个栈
%  
\chapter{调度}
\label{CH:SCHED}

任何操作系统都可能运行比计算机 CPU 数量更多的进程，因此需要一种计划来在进程之间对 CPU 进行时间共享。理想情况下，这种共享对用户进程应该是透明的。一种常见的方法是通过 \indextext{多路复用} 为每个进程提供其拥有虚拟 CPU 的假象，即将进程多路复用到硬件 CPU 上。本章解释 xv6 如何实现这种多路复用。

在继续阅读本章之前，请阅读 {\tt kernel/proc.h}、{\tt kernel/swtch.S} 以及 {\tt kernel/proc.c} 中的 {\tt yield()}、{\tt sched()} 和 {\tt schedule()}。

%% 
\section{多路复用}
%% 

Xv6 通过两种情况下在每个 CPU 上从一个进程切换到另一个进程来实现多路复用。第一种情况是，当进程执行会阻塞（必须等待）的系统调用时，例如 \lstinline{read} 或 \lstinline{wait}，xv6 会进行切换。第二种情况是，xv6 会定期强制进行切换，以应对那些长时间进行计算而不阻塞的进程。前者称为自愿切换，后者称为非自愿切换。

实现多路复用带来了一些挑战。第一，如何从一个进程切换到另一个进程？基本思想是保存和恢复 CPU 寄存器，尽管这一过程无法用 C 语言表达，因此实现起来比较棘手。第二，如何以对用户进程透明的方式强制进行切换？Xv6 使用了标准技术，即利用硬件定时器的中断来驱动上下文切换。第三，所有 CPU 都在同一组进程之间切换，因此需要一个锁计划来避免错误，例如两个 CPU 决定同时运行同一个进程。第四，当进程退出时，其内存和其他资源必须被释放，但它无法自己完成所有这些工作。第五，多核机器的每个 CPU 必须记住它正在执行哪个进程，以便系统调用影响正确进程的内核状态。

%% 
\section{上下文切换概述}
%% 

术语“上下文切换”指的是 CPU 停止执行一个内核线程（通常是为了之后恢复执行），并开始执行另一个不同内核线程所涉及的步骤；这种切换是多路复用的核心。Xv6 并不直接从一个进程的内核线程上下文切换到另一个进程的内核线程；相反，一个内核线程通过上下文切换到该 CPU 的“调度器线程”来放弃 CPU，然后调度器线程选择另一个不同进程的内核线程来运行，并上下文切换到该线程。

\begin{figure}[t]
\center
\includegraphics[scale=0.5]{fig/switch.pdf}
\caption{从一个用户进程切换到另一个用户进程。在此示例中，xv6 运行在一个 CPU 上（因此只有一个调度器线程）。}
\label{fig:switch}
\end{figure}

从更广的 scope 来看，从一个用户进程切换到另一个用户进程所涉及的步骤如图~\ref{fig:switch} 所示：从旧进程的用户空间陷入（系统调用或中断）到其内核线程，然后上下文切换到当前 CPU 的调度器线程，再上下文切换到新进程的内核线程，最后通过陷阱返回到用户级进程。

%% 
\section{代码：上下文切换}
%% 

{\tt kernel/swtch.S} 中的函数 {\tt swtch()} 包含了线程上下文切换的核心：它保存被切换出去的线程的 CPU 寄存器，并恢复之前保存的将要切换到的线程的寄存器。这样做之所以足够，其基本原因是线程的状态由内存中的数据（例如其栈）加上其 CPU 寄存器组成；线程内存不需要保存和恢复，因为不同的线程将其数据保存在 RAM 的不同区域；但 CPU 只有一组寄存器，因此它们必须在线程之间进行切换（保存和恢复）。

每个线程的 {\tt struct proc} 包含一个 {\tt struct context}，用于在该线程未运行时保存其寄存器的值。CPU 的调度器线程的 {\tt struct context} 保存在该 CPU 的 {\tt struct cpu} 中。当线程 X 希望切换到线程 Y 时，线程 X 调用 {\tt swtch(\&X's context, \&Y's context)}。{\tt swtch()} 将当前 CPU 寄存器保存到 X 的上下文中，然后将 Y 的上下文内容加载到 CPU 寄存器中，然后返回。

以下是 {\tt swtch} 的简化副本：

\begin{verbatim}
swtch:
        sd ra, 0(a0)
        sd sp, 8(a0)
        sd s0, 16(a0)
        ...
        sd s11, 104(a0)

        ld ra, 0(a1)
        ld sp, 8(a1)
        ld s0, 16(a1)
        ...
        ld s11, 104(a1)
        
        ret
\end{verbatim}

{\tt a0} 保存第一个函数参数，{\tt a1} 保存第二个参数；在此例中，即两个 {\tt struct context} 指针。{\tt 16(a0)} 指的是从 {\tt a0} 指向的内存地址偏移 16 字节处；参考 {\tt kernel/proc.h} 中 {\tt struct context} 的定义 \lineref{kernel/proc.h:/^struct.context/}，这是名为 {\tt s0} 的结构体字段。

{\tt swtch} 的 {\tt ret} 返回到哪里？它返回到 {\tt ra} 寄存器指向的指令。在线程 X 调用 {\tt swtch()} 切换到 Y 的例子中，当 {\tt ret} 执行时，{\tt ra} 刚刚从 Y 的 {\tt struct context} 中加载进来。而 Y 的 {\tt struct context} 中的 {\tt ra} 最初是在 Y 过去放弃 CPU 时，由 Y 调用 {\tt swtch} 保存的。因此，{\tt ret} 返回到 Y 调用 {\tt swtch()} 之后的那个指令处；也就是说，X 对 {\tt swtch()} 的调用就好像是从 Y 原来的 {\tt swtch()} 调用返回一样。并且 {\tt sp} 将是 Y 的栈指针，因为 {\tt swtch} 从 Y 的 {\tt struct context} 中加载了 {\tt sp}；因此在返回时，Y 将在自己的栈上执行。{\tt swtch()} 不需要直接保存或恢复程序计数器；保存和恢复 {\tt ra} 就足够了。

\lstinline{swtch}
\lineref{kernel/swtch.S:/swtch/}
保存被调用者保存的寄存器（{\tt ra,sp,s0..s11}），但不保存调用者保存的寄存器。RISC-V 调用约定要求，如果代码需要在函数调用过程中保留调用者保存寄存器中的值，编译器必须生成指令在函数调用前将该寄存器保存到栈上，并在函数返回时从栈上恢复。因此 \lstinline{swtch} 可以依赖调用它的函数已经保存了调用者保存的寄存器（或者调用函数根本不需要这些寄存器中的值）。

%% 
\section{代码：调度}
%% 

% 为什么需要一个单独的线程

上一节研究了 \indexcode{swtch} 的内部细节；现在让我们把 \lstinline{swtch} 当作一个已知函数，来看看如何从一个进程的内核线程通过调度器切换到另一个进程。调度器以每个 CPU 一个特殊线程的形式存在，每个线程都运行 \lstinline{scheduler} 函数。该函数负责选择接下来运行哪个进程。每个 CPU 都有自己的调度器线程，因为可能同时有多个 CPU 在寻找可运行的进程。进程切换总是通过调度器线程，而不是直接从一进程到另一进程，这是为了避免某些情况下没有栈可以用来执行调度器（例如，如果旧进程已经退出，或者当前没有其他进程想要运行）。

一个想要放弃 CPU 的进程必须获取自己的进程锁 \indexcode{p->lock}，释放它持有的任何其他锁，更新自己的状态 (\lstinline{p->state})，然后调用 \indexcode{sched}。你可以在 \lstinline{yield} \lineref{kernel/proc.c:/^yield/}、\texttt{sleep} 和 \texttt{kexit} 中看到这个序列。\indexcode{sched} 调用 \indexcode{swtch} 将当前上下文保存到 \lstinline{p->context} 中，并切换到 \indexcode{cpu->context} 中的调度器上下文。\lstinline{swtch} 在调度器的栈上返回，就好像 \indexcode{scheduler} 的 \lstinline{swtch} 已经返回了一样 \lineref{kernel/proc.c:/swtch.*contex.*contex/}。

\lstinline{scheduler} \lineref{kernel/proc.c:/^scheduler/} 运行一个循环：找到一个要运行的进程，通过 {\tt swtch()} 切换到它，最终它会通过 {\tt swtch()} 切换回调度器，调度器继续循环。调度器遍历进程表，寻找一个可运行的进程，即 \lstinline{p->state} \lstinline{==} \lstinline{RUNNABLE} 的进程。一旦找到这样的进程，它就设置每 CPU 的当前进程变量 \lstinline{c->proc}，将该进程标记为 \lstinline{RUNNING}，然后调用 \indexcode{swtch} 开始运行它 \linerefs{kernel/proc.c:/Switch.to/,/swtch/}。在过去的某个时刻，目标进程一定调用过 {\tt swtch()}；调度器对 {\tt swtch()} 的调用实际上就是从那个更早的调用返回。图~\ref{fig:inout} 说明了这种模式。

\begin{figure}[t]
\input{fig/switch.tex}
\caption{{\tt swtch()} 总是以调度器线程作为源或目标，并且相关的 {\tt p->lock} 始终被持有。}
\label{fig:inout}
\end{figure}

xv6 在调用 \lstinline{swtch} 的过程中持有 \indexcode{p->lock}：调用 \indexcode{swtch} 的线程获取锁，但锁是在目标线程中 \lstinline{swtch} 返回后才释放的。这种安排不同寻常：更常见的是获取锁的线程也负责释放锁。Xv6 的上下文切换打破了这一惯例，因为 \indexcode{p->state} 和 \lstinline{p->context} 必须以原子方式一起更新。例如，如果在调用 \indexcode{swtch} 之前释放了 \lstinline{p->lock}，那么另一个 CPU $c$ 可能会因为该进程的状态是 \lstinline{RUNNABLE} 而决定运行它。CPU $c$ 将调用 \lstinline{swtch}，并从 \lstinline{p->context} 中恢复状态，而此时原来的 CPU 仍在向 \lstinline{p->context} 保存状态。其结果将是该进程在 CPU $c$ 上以部分保存的寄存器状态恢复，并且两个 CPU 将使用同一个栈，这将导致混乱。一旦 \lstinline{yield} 开始修改一个正在运行的进程的状态，使其变成 \lstinline{RUNNABLE}，就必须一直持有 \lstinline{p->lock}，直到该进程已保存其所有寄存器并且调度器在其栈上运行。最早的正确释放点是在 \lstinline{scheduler}（在其自己的栈上运行）清除 \lstinline{c->proc} 之后。类似地，一旦 \lstinline{scheduler} 开始将一个 \lstinline{RUNNABLE} 的进程转换为 \lstinline{RUNNING}，就不能释放锁，直到该进程的内核线程完全运行起来（即在 \lstinline{swtch} 之后，例如在 \lstinline{yield} 中）。

有一种情况是调度器对 \indexcode{swtch} 的调用不会最终进入 \indexcode{sched}。\lstinline{allocproc} 将新进程的上下文 \lstinline{ra} 寄存器设置为 \indexcode{forkret} \lineref{kernel/proc.c:/^forkret/}，这样它的第一次 \lstinline{swtch} 就会“返回”到该函数的开头。\lstinline{forkret} 的存在是为了释放 \indexcode{p->lock} 并设置一些控制寄存器和陷阱帧字段，这些是返回用户空间所必需的。最后，{\tt forkret} 模拟从系统调用返回到用户空间的正常返回路径。

%% 
\section{代码：mycpu 和 myproc}
%% 

Xv6 经常需要指向当前进程的 \lstinline{proc} 结构的指针。在单处理器上，可以使用一个全局变量指向当前的 \lstinline{proc}。但在多核机器上这行不通，因为每个 CPU 执行的是不同的进程。解决这个问题的方法是利用每个 CPU 都拥有自己一组寄存器这一事实。

当某个 CPU 在内核中执行时，xv6 确保该 CPU 的 \lstinline{tp} 寄存器始终保存着该 CPU 的 hartid。RISC-V 为其每个 CPU 编号，并赋予每个 CPU 一个唯一的 \indextext{hartid}。\indexcode{mycpu} \lineref{kernel/proc.c:/^mycpu/} 使用 \lstinline{tp} 作为 \lstinline{cpu} 结构数组的索引，并返回当前 CPU 对应的那个结构。一个 \indexcode{struct cpu} \lineref{kernel/proc.h:/^struct.cpu/} 持有一个指向当前在该 CPU 上运行的进程（如果有的话）的 \lstinline{proc} 结构指针、该 CPU 调度器线程的保存寄存器，以及用于管理中断禁用的嵌套自旋锁计数。

确保 CPU 的 \lstinline{tp} 保存着该 CPU 的 hartid 有点复杂，因为用户代码可以随意修改 \lstinline{tp}。\lstinline{start} 在 CPU 启动序列的早期，仍处于机器模式时，就设置了 \lstinline{tp} 寄存器 \lineref{kernel/start.c:/w_tp/}。在准备返回用户空间时，\lstinline{prepare_return} 将 \lstinline{tp} 保存到跳板页中，以防用户代码修改它。最后，\lstinline{uservec} 在从用户空间进入内核时，恢复那个保存的 \lstinline{tp} \lineref{kernel/trampoline.S:/make tp hold/}。编译器保证在内核代码中永远不会修改 \lstinline{tp}。如果 xv6 能在需要时直接向 RISC-V 硬件询问当前的 hartid 会更方便，但 RISC-V 只允许在机器模式下这样做，而不允许在监管者模式下进行。

\lstinline{cpuid} 和 \lstinline{mycpu} 的返回值是不可靠的：如果定时器中断导致线程让出 CPU，并在之后在不同的 CPU 上恢复执行，那么先前返回的值将不再正确。为了避免这个问题，xv6 要求代码在调用 {\tt cpuid()} 或 {\tt mycpu()} 之前禁用中断，并且只有在使用完返回值之后才能重新启用中断。

函数 \indexcode{myproc} \lineref{kernel/proc.c:/^myproc/} 返回当前 CPU 上正在运行的进程的 \lstinline{struct proc} 指针。\lstinline{myproc} 会禁用中断，调用 \lstinline{mycpu}，从 \lstinline{struct cpu} 中取出当前进程指针 (\lstinline{c->proc})，然后重新启用中断。即使中断是启用的，\lstinline{myproc} 的返回值也可以安全使用：如果定时器中断将调用进程移到了另一个不同的 CPU 上，它的 \lstinline{struct proc} 指针将保持不变。

%% 
\section{现实世界}
%% 

Xv6 调度器实现了一个简单的调度策略，轮流运行每个进程。这种策略称为 \indextext{轮转调度}。真正的操作系统实现了更复杂的策略，例如允许进程拥有优先级。其思想是，调度器会优先选择可运行的高优先级进程，而不是可运行的低优先级进程。这些策略可能会变得复杂，因为通常存在相互竞争的目标：例如，操作系统可能还希望保证公平性和高吞吐量。

%% 
\section{练习}
%% 

\begin{enumerate}

\item 修改 xv6，使得从一个进程的内核线程切换到另一个进程时只使用一次上下文切换，而不是通过调度器线程进行切换。让出 CPU 的线程需要自己选择下一个线程并调用 \texttt{swtch}。挑战在于：防止多个 CPU 意外执行同一个线程；正确使用锁；以及避免死锁。

\end{enumerate}
